% ========================================================
% BÀI TẬP TOÁN 5 - KẾT NỐI TRI THỨC VỚI CUỘC SỐNG
% Bài 12: Viết số đo đại lượng dưới dạng số thập phân
% Định dạng: BTTeX Part III - Trắc nghiệm điền số ngắn
% ========================================================

% Câu 1
\begin{ex}
  Viết số đo $3\text{ m } 5\text{ dm}$ dưới dạng số thập phân có đơn vị mét là $3{,}a\text{ m}$. Chữ số $a$ là bao nhiêu?
  \shortans{5}
  \loigiai{
    $3\text{ m } 5\text{ dm} = 3\frac{5}{10}\text{ m} = 3{,}5\text{ m}$. Chữ số $a$ là 5.
  }
\end{ex}

% Câu 2
\begin{ex}
  Viết số đo $4\text{ m } 25\text{ cm}$ dưới dạng số thập phân có đơn vị mét là $4{,}bc\text{ m}$. Tổng hai chữ số $b + c$ là bao nhiêu?
  \shortans{7}
  \loigiai{
    $4\text{ m } 25\text{ cm} = 4\frac{25}{100}\text{ m} = 4{,}25\text{ m}$. Tổng là $2 + 5 = 7$.
  }
\end{ex}

% Câu 3
\begin{ex}
  Viết số đo $8\text{ cm } 6\text{ mm}$ dưới dạng số thập phân có đơn vị xăng-ti-mét là $8{,}6\text{ cm}$. Chữ số hàng phần mười là bao nhiêu?
  \shortans{6}
  \loigiai{
    $8\text{ cm } 6\text{ mm} = 8{,}6\text{ cm}$. Chữ số hàng phần mười là 6.
  }
\end{ex}

% Câu 4
\begin{ex}
  Viết số đo $2\text{ km } 350\text{ m}$ dưới dạng số thập phân có đơn vị ki-lô-mét là $2{,}35\text{ km}$. Chữ số hàng phần trăm là bao nhiêu?
  \shortans{5}
  \loigiai{
    $2\text{ km } 350\text{ m} = 2{,}35\text{ km}$. Chữ số hàng phần trăm là 5.
  }
\end{ex}

% Câu 5
\begin{ex}
  Viết số đo $7\text{ km } 45\text{ m}$ dưới dạng số thập phân có đơn vị ki-lô-mét là $7{,}045\text{ km}$. Chữ số hàng phần mười của số này là bao nhiêu?
  \shortans{0}
  \loigiai{
    $7\text{ km } 45\text{ m} = 7\frac{45}{1000}\text{ km} = 7{,}045\text{ km}$. Chữ số hàng phần mười là 0.
  }
\end{ex}

% Câu 6
\begin{ex}
  Viết số đo $5\text{ tấn } 8\text{ tạ}$ dưới dạng số thập phân có đơn vị tấn là $5{,}8\text{ tấn}$. Chữ số hàng phần mười là bao nhiêu?
  \shortans{8}
  \loigiai{
    $5\text{ tấn } 8\text{ tạ} = 5{,}8\text{ tấn}$. Chữ số hàng phần mười là 8.
  }
\end{ex}

% Câu 7
\begin{ex}
  Viết số đo $3\text{ tấn } 250\text{ kg}$ dưới dạng số thập phân có đơn vị tấn là $3{,}25\text{ tấn}$. Chữ số hàng phần mười là bao nhiêu?
  \shortans{2}
  \loigiai{
    $3\text{ tấn } 250\text{ kg} = 3{,}25\text{ tấn}$. Chữ số hàng phần mười là 2.
  }
\end{ex}

% Câu 8
\begin{ex}
  Viết số đo $6\text{ tấn } 50\text{ kg}$ dưới dạng số thập phân có đơn vị tấn là $6{,}05\text{ tấn}$. Chữ số hàng phần trăm là bao nhiêu?
  \shortans{5}
  \loigiai{
    $6\text{ tấn } 50\text{ kg} = 6\frac{50}{1000}\text{ tấn} = 6{,}05\text{ tấn}$. Chữ số hàng phần trăm là 5.
  }
\end{ex}

% Câu 9
\begin{ex}
  Viết số đo $4\text{ kg } 75\text{ g}$ dưới dạng số thập phân có đơn vị ki-lô-gam là $4{,}075\text{ kg}$. Chữ số hàng phần nghìn là bao nhiêu?
  \shortans{5}
  \loigiai{
    $4\text{ kg } 75\text{ g} = 4{,}075\text{ kg}$. Chữ số hàng phần nghìn là 5.
  }
\end{ex}

% Câu 10
\begin{ex}
  Viết số đo $12\text{ kg } 8\text{ g}$ dưới dạng số thập phân có đơn vị ki-lô-gam là $12{,}008\text{ kg}$. Chữ số hàng phần trăm là bao nhiêu?
  \shortans{0}
  \loigiai{
    $12\text{ kg } 8\text{ g} = 12{,}008\text{ kg}$. Chữ số hàng phần trăm là 0.
  }
\end{ex}

% Câu 11
\begin{ex}
  Viết số đo $5\text{ m}^2\text{ } 35\text{ dm}^2$ dưới dạng số thập phân có đơn vị mét vuông là $5{,}35\text{ m}^2$. Chữ số hàng phần mười là bao nhiêu?
  \shortans{3}
  \loigiai{
    $5\text{ m}^2\text{ } 35\text{ dm}^2 = 5{,}35\text{ m}^2$. Chữ số hàng phần mười là 3.
  }
\end{ex}

% Câu 12
\begin{ex}
  Viết số đo $8\text{ m}^2\text{ } 7\text{ dm}^2$ dưới dạng số thập phân có đơn vị mét vuông là $8{,}07\text{ m}^2$. Chữ số hàng phần trăm là bao nhiêu?
  \shortans{7}
  \loigiai{
    $8\text{ m}^2\text{ } 7\text{ dm}^2 = 8\frac{7}{100}\text{ m}^2 = 8{,}07\text{ m}^2$. Chữ số hàng phần trăm là 7.
  }
\end{ex}

% Câu 13
\begin{ex}
  Viết số đo $15\text{ dm}^2\text{ } 8\text{ cm}^2$ dưới dạng số thập phân có đơn vị đề-xi-mét vuông là $15{,}08\text{ dm}^2$. Chữ số hàng phần mười là bao nhiêu?
  \shortans{0}
  \loigiai{
    $15\text{ dm}^2\text{ } 8\text{ cm}^2 = 15{,}08\text{ dm}^2$. Chữ số hàng phần mười là 0.
  }
\end{ex}

% Câu 14
\begin{ex}
  Đổi $4{,}5\text{ m}$ thành bao nhiêu đề-xi-mét?
  \shortans{45}
  \loigiai{
    $4{,}5\text{ m} = 4{,}5 \times 10 = 45\text{ dm}$.
  }
\end{ex}

% Câu 15
\begin{ex}
  Đổi $2{,}75\text{ m}$ thành bao nhiêu xăng-ti-mét?
  \shortans{275}
  \loigiai{
    $2{,}75\text{ m} = 2{,}75 \times 100 = 275\text{ cm}$.
  }
\end{ex}

% Câu 16
\begin{ex}
  Đổi $1{,}5\text{ km}$ thành bao nhiêu mét?
  \shortans{1500}
  \loigiai{
    $1{,}5\text{ km} = 1{,}5 \times 1\,000 = 1\,500\text{ m}$.
  }
\end{ex}

% Câu 17
\begin{ex}
  Đổi $0{,}75\text{ tấn}$ thành bao nhiêu ki-lô-gam?
  \shortans{750}
  \loigiai{
    $0{,}75\text{ tấn} = 0{,}75 \times 1\,000 = 750\text{ kg}$.
  }
\end{ex}

% Câu 18
\begin{ex}
  Đổi $3{,}2\text{ tạ}$ thành bao nhiêu yến?
  \shortans{32}
  \loigiai{
    $3{,}2\text{ tạ} = 3{,}2 \times 10 = 32\text{ yến}$.
  }
\end{ex}

% Câu 19
\begin{ex}
  Đổi $2{,}5\text{ m}^2$ thành bao nhiêu đề-xi-mét vuông?
  \shortans{250}
  \loigiai{
    $2{,}5\text{ m}^2 = 2{,}5 \times 100 = 250\text{ dm}^2$.
  }
\end{ex}

% Câu 20
\begin{ex}
  Đổi $0{,}08\text{ m}^2$ thành bao nhiêu xăng-ti-mét vuông?
  \shortans{800}
  \loigiai{
    $0{,}08\text{ m}^2 = 0{,}08 \times 10\,000 = 800\text{ cm}^2$.
  }
\end{ex}

% Câu 21
\begin{ex}
  Một sợi dây dài $3\text{ m } 60\text{ cm}$ được cắt thành các đoạn dài $0{,}9\text{ m}$. Cắt được bao nhiêu đoạn?
  \shortans{4}
  \loigiai{
    $3\text{ m } 60\text{ cm} = 3{,}6\text{ m}$. Số đoạn là: $3{,}6 : 0{,}9 = 4$ (đoạn).
  }
\end{ex}

% Câu 22
\begin{ex}
  Một con bò nặng $3\text{ tạ } 50\text{ kg}$. Viết khối lượng con bò dưới dạng số thập phân có đơn vị tấn là $0{,}35\text{ tấn}$. Chữ số hàng phần trăm là bao nhiêu?
  \shortans{5}
  \loigiai{
    $3\text{ tạ } 50\text{ kg} = 350\text{ kg} = 0{,}35\text{ tấn}$. Chữ số hàng phần trăm là 5.
  }
\end{ex}

% Câu 23
\begin{ex}
  Một khu đất hình chữ nhật có chiều dài $20\text{ m}$, chiều rộng $15\text{ m}$. Diện tích khu đất viết dưới dạng số thập phân có đơn vị héc-ta là $0{,}03\text{ ha}$. Chữ số hàng phần trăm là bao nhiêu?
  \shortans{3}
  \loigiai{
    Diện tích là $20 \times 15 = 300\text{ m}^2 = \frac{300}{10000}\text{ ha} = 0{,}03\text{ ha}$. Chữ số hàng phần trăm là 3.
  }
\end{ex}

% Câu 24
\begin{ex}
  Viết số đo $500\text{ m}$ dưới dạng số thập phân có đơn vị ki-lô-mét là $0{,}5\text{ km}$. Phần nguyên của số đó là bao nhiêu?
  \shortans{0}
  \loigiai{
    Phần nguyên là 0.
  }
\end{ex}

% Câu 25
\begin{ex}
  Một mảnh vườn có chu vi $48\text{ m}$. Nửa chu vi mảnh vườn đổi ra ki-lô-mét là $0{,}024\text{ km}$. Chữ số hàng phần nghìn là bao nhiêu?
  \shortans{4}
  \loigiai{
    Nửa chu vi là $24\text{ m} = 0{,}024\text{ km}$. Chữ số hàng phần nghìn là 4.
  }
\end{ex}
